% \paragraph{Conclude}
% Summarize the work you did. Show how well you reached the goal you
% specified in the introduction.
%
% \paragraph{Outlook}
% Which new questions arise based on your results? If you would continue
% to work on this topic, what would your next work item be?
\chapter{Conclusions and outlook}\label{ch:conclusion}

This thesis set out to design a packet buffer that bridges the mismatch between packets
arriving individually and interleaved across many flows, and a forward error correction
encoder that must instead consume them per-flow and in fixed-size generations. The proposed
design follows a shared-memory architecture: packet payloads are held in a single pool of
HBM-backed memory shared across all flows, while a lightweight, per-flow queue of packet
addresses (rather than payloads) tracks each flow's pending packets and exposes them to the
encoder in complete generations. \\

The design was evaluated through behavioural simulation, using a testbench built around a
configurable HBM memory model, a parametrizable UDP packet generator, and a simple FEC encoder
model, all communicating over AXI-Stream. Across the three sweeps performed in
chapter \ref{ch:evaluation}, the design met the requirements set out in
section \ref{system_requirements} for the example configuration used throughout this thesis
(generation size 32, 128 flows, a free memory queue depth of 1024, and a per-flow queue depth
of 64): throughput stayed close to the 100 Gb/s target regardless of packet size, and the
minimum packet latency attributable to the buffer itself stayed under the 10~\textmu s target. \\

As discussed in chapter \ref{ch:discussion}, packet latency is shown
to depend heavily on how packets from different flows happen to be interleaved rather than on
the buffer's own architecture, and a similar, workload-driven limit was identified on
throughput for very small packets, arising from the one-transaction-per-packet ingest
interface rather than from the buffer's internal pipeline. \\

As it is already covered in chapter \ref{ch:discussion}, 
the directions for addressing the aforementioned points would be 
consolidating the 128 per-flow packet descriptor queues into a single BRAM-based structure,
adding the slave configuration registers already contemplated by the Open NIC project to allow
runtime and per-flow configuration, and replacing the current tuple lookup with a proper hash
table; all of which can easily be fitted into the current design without requiring a substantial re-factor or re-design. \\

Lastly, it is crucial to note that besides the improvements on the current design, in order to 
actually have full FEC system the RX side of the system must also be implemented. While the 
details on how to do it are outside the scope of this thesis, it is worth emphasizing that the 
problem is highly symmetrical, and that the large majority of the modules could be reused for 
that purpose, specially the ones related to the management of the HBM pool. \\

Overall, this thesis shows that a shared-memory, address-queue-based packet buffer is a viable
way to reconcile per-packet arrival with per-generation, per-flow FEC encoding on an
HBM-equipped SmartNIC platform, and provides both a working baseline implementation and a clear
set of next steps for improving the set of features for a future SmartNIC transmit pipeline.
